\section{Related Work}

The closest prior work is Yim et al.~\cite{yim2023entanglement}, which detects leg entanglement from proprioception and, when a leg snags, executes a disentanglement maneuver as that leg advances through its swing phase, succeeding in 14 of 16 laboratory trials and in outdoor tests. It establishes entanglement as a real and tractable failure mode, but it also leaves clear gaps. Detection is largely rule-based and reactive, tied to gait-phase timing rather than to a learned temporal model; the output is effectively a binary trigger, with no attribution of which leg is caught and no measure of how strongly it is held; and sensing is fused with the reactive control loop rather than exposed as a reusable estimate. We keep the proprioceptive premise but address these gaps: we integrate evidence over a fixed $0.40$~s window with a learned causal model instead of an instantaneous rule, we emit structured outputs (which leg, and a calibrated per-leg severity in $[0,1]$), and we treat detection as a self-contained perception problem that publishes a state estimate, decoupled from any actuation policy. The contrast is in what is inferred and how it is reported, not in closing a control loop.

Proprioceptive learning is otherwise well established for legged systems, from locomotion policies that transfer to hardware~\cite{hwangbo2019learning,lee2020learning} to probabilistic contact and event estimation~\cite{camurri2017contact,bledt2018contact}. Entanglement differs in kind: it is an externally imposed kinematic constraint, applied at an unpredictable time and place, that must be localized to a leg and graded by how strongly it resists, rather than classified as contact versus no contact. For the model we build on temporal convolutional networks~\cite{bai2018tcn} and the dilated causal convolutions of WaveNet~\cite{oord2016wavenet}, whose causality and stride-spanning receptive field suit an at-every-timestep detector that must run online. Finally, because a detector is only usable if its confidence maps to a trustworthy threshold, we apply temperature scaling~\cite{guo2017calibration,niculescu2005predicting}, mainly to de-saturate probabilities that would otherwise pin near $1.0$.
